\makeatletter
\renewcommand {\@degree@string} {Master of Science}
\makeatother

\title{Auto-Adapter: Robot-Specific Driver Synthesis with Independent Closed-Loop Validation}
\author{Yifan Wang}
\department{Department of Computer Science}

\maketitle
\makedeclaration

\begin{abstract}
Deploying a tool-using large language model (LLM) agent on a new robot requires
a tool application programming interface (API) and a robot-specific driver
behind it. Existing agent frameworks generally assume that this software layer
has already been written. This thesis asks whether an LLM can instead design
the API, define how it should be tested, and implement the driver itself.
Auto-Adapter reads a robot description and a collection of target tasks,
proposes a small set of reusable robot operations, and generates the code that
exposes them to an agent. The proposed operations are evaluated against
measurable criteria in simulation, using observations and verdicts that remain
separate from the generated code. When a driver fails, Auto-Adapter permits a
fixed number of revision attempts. The study therefore distinguishes between
designing robot tools, implementing them, validating their behaviour, and
using them to complete tasks.

The evaluation comprises three experiments.
\textbf{Experiment 1---Auto-Adapter Framework.} The first experiment evaluates
the complete framework, from designing reusable robot operations to generating
and independently validating a robot-specific driver. Its purpose is to
determine whether Auto-Adapter can construct the robot-facing software required
by a tool-using agent.

\textbf{Experiment 2---AutoAdapter-Bench.} The benchmark evaluates LLMs
separately as creators and users of robot tools. Its purpose is to compare these
abilities and establish whether successful tool use also indicates successful
tool synthesis.

\textbf{Experiment 3---Cross-Robot Evaluation.} The final experiment applies
Auto-Adapter across different robot configurations. Its purpose is to analyse
how observed performance differs between robots and identify where the
framework succeeds or fails.

[RESULT NEEDED: report the principal Experiment~1 synthesis and resource
outcomes, the Experiment~2 component-analysis, capability-design, and
validation-design outcomes, and the matched Experiment~3 Experience effect.]
The findings apply only to the
declared model endpoints, tasks, robot configurations, and Direct-MuJoCo
conditions; they do not establish real-SDK fidelity, hardware safety,
sim-to-real transfer, universal robot or model superiority, autonomous
self-improvement, or generalisation beyond the evaluated conditions.

Within the evaluated scope, the thesis makes three scientific contributions:
\vspace{-.65\baselineskip}
\begin{enumerate}[leftmargin=*,labelsep=.8em,nosep]

  \item A layer-separated benchmark of LLMs as producers and users of robot
  software. B1 evaluates robot-specific driver synthesis under fixed
  capability interfaces and private validation, whereas B2 evaluates
  controller backbones against fixed, validated
  driver--interface--adapter stacks. Their factors, statistical units,
  outcomes, and verdicts remain separate, so downstream task success is not
  counted as evidence of driver synthesis.

  \item An evidence-governed synthesis lifecycle for model-authored robot
  capabilities and controlled adaptation. Task-Grounded Capability Design
  derives the capability layer and its acceptance requirements from
  source-backed tasks; Driver Synthesis and bounded same-run Repair implement
  that layer. Terminal Evolution may propose a restricted Experience record,
  but only human-reviewed Experience can enter a later matched run.

  \item An implementation-blind, physics-grounded test-oracle architecture for
  generated robot software. The Independent Validation Compiler audits source
  preservation and compiles accepted criteria without inspecting the candidate
  implementation. The trusted Harness supplies Direct-MuJoCo observations,
  enforces physical-integrity guards, and owns final verdicts, thereby
  separating criterion compilation and structural conformance from
  actuator-mediated behaviour.

\end{enumerate}

Within the declared Direct-MuJoCo scope, the thesis also makes three bounded,
industry-relevant contributions:
\vspace{-.65\baselineskip}
\begin{enumerate}[leftmargin=*,labelsep=.8em,nosep]

  \item A bounded and inspectable workflow for LLM-assisted robot integration.
  Before each controlled comparison, Auto-Adapter fixes the admissible
  per-robot inputs and callable contract, bounds generation and Repair, records
  model and resource use, and assigns private physics-grounded evaluation to
  the trusted Harness. Candidate revisions, model calls, token use, cost, wall
  time, observations, and acceptance decisions can therefore be inspected
  within the declared workflow.

  \item A reusable experimental substrate for embodied-AI research and
  development. Each admitted experimental stack exposes a stable capability
  interface backed by a driver that has passed independent Direct-MuJoCo
  capability validation. High-level-controller backbones can therefore be
  compared against the same robot-facing software boundary without changing
  the underlying integration within that comparison.

  \item A resource-accounted basis for role-specific LLM selection. B1
  supports conditional comparison of Producer backbones and trusted-skeleton
  access for driver synthesis, whereas B2 supports conditional comparison of
  controller backbones after the robot-facing stack has been fixed. Reporting
  behavioural outcomes with calls, tokens, cost, wall time, Repair, and failure
  records where applicable supports role-specific trade-off analysis rather
  than a universal model ranking.

\end{enumerate}

\end{abstract}

\begin{acknowledgements}
I would like to express my sincere gratitude to my supervisors,
Professor Philip Treleaven, Dr Adriano Koshiyama, Mr Zekun Wu, and
Mr Kleyton Costa, for their guidance, thoughtful feedback, and encouragement
throughout this project. Their questions helped me refine the research problem,
strengthen the experimental design, and communicate the scope and limitations
of the work more clearly.

I am also grateful to the staff and students in UCL's Department of Computer
Science for the discussions and practical support that contributed to this
research. I also thank the maintainers of the open-source software, simulation
models, and robot assets on which the experimental system was built.

Finally, I would like to thank my family and friends for their patience,
understanding, and constant support throughout the MSc. Their encouragement
made the most demanding stages of this project much easier to navigate.
\end{acknowledgements}

\setcounter{tocdepth}{2}
% Setting this higher means you get contents entries for
%  more minor section headers.

\tableofcontents
\listoffigures
\listoftables
